% GENERATED FILE. Edit canonical lessons, then run npm run generate:books.
% canonical-source-hash: fnv1a64:5791909a51be4a59
% canonical-lessons: GE-C191-vorbereiten, GE-C191-waehlen, GE-C191-wecken, GE-C191-wiegen, GE-C191-zeichnen, GE-R191-first-pass-words-from-chapters-183-186, GE-R191-second-pass-words-from-chapters-187-191

\chapter{To prepare, To choose, To wake, To weigh, To draw}
\label{ch:ge-to-prepare-to-choose-to-wake-to-weigh-to-draw}
\hlchaptermodality{191}

\begin{chapteropening}
\textbf{By the end of this chapter:} \emph{I can say to prepare, to choose, to wake, to weigh and to draw.}

Complete the last lesson of chapter 191: I can say to prepare, to choose, to wake, to weigh and to draw.
\end{chapteropening}

\section[vorbereiten]{vorbereiten — to prepare}

\label{lesson:GE-C191-vorbereiten}



\begin{quote}
Before the new one: say the German for to postpone, then the German for to use.
\end{quote}

\subsection*{You'll want to know: vorbereiten}
\textbf{vorbereiten} — \textquotedblleft{}to prepare\textquotedblright{}.

\begin{grammarlens}[title={What you've built}]
One more word for this chapter.
\end{grammarlens}

\subsection*{Your turn}
\begin{itemize}
\raggedright
  \item \emph{Say it:} \emph{vorbereiten}
  \item \emph{Say it:} \emph{vorbereiten}, once more
  \item \emph{Say it:} say \emph{vorbereiten}
  \item \emph{Recall:} say the German for to postpone, then the German for to use, then say \emph{vorbereiten} again
\end{itemize}

\begin{tcolorbox}[breakable,colback=teal!4,colframe=teal!35!black,title={Before you move on}]
What does vorbereiten mean? (\textquotedblleft{}To prepare\textquotedblright{}.) Say it once more.
\end{tcolorbox}
\section[wählen]{wählen — to choose, to vote, to dial}

\label{lesson:GE-C191-waehlen}



\begin{quote}
Before the new one: say the German for to use, then the German for to prepare.
\end{quote}

\subsection*{You'll want to know: wählen}
\textbf{wählen} — \textquotedblleft{}to choose, to vote, to dial\textquotedblright{}.

\begin{grammarlens}[title={What you've built}]
One more word for this chapter.
\end{grammarlens}

\subsection*{Your turn}
\begin{itemize}
\raggedright
  \item \emph{Say it:} \emph{wählen}
  \item \emph{Say it:} \emph{wählen}, once more
  \item \emph{Say it:} say \emph{wählen}
  \item \emph{Recall:} say the German for to use, then the German for to prepare, then say \emph{wählen} again
\end{itemize}

\begin{tcolorbox}[breakable,colback=teal!4,colframe=teal!35!black,title={Before you move on}]
What does wählen mean? (\textquotedblleft{}To choose, to vote, to dial\textquotedblright{}.) Say it once more.
\end{tcolorbox}
\section[wecken]{wecken — to wake (someone)}

\label{lesson:GE-C191-wecken}



\begin{quote}
Before the new one: say the German for to prepare, then the German for to choose.
\end{quote}

\subsection*{You'll want to know: wecken}
\textbf{wecken} — \textquotedblleft{}to wake (someone)\textquotedblright{}.

\begin{grammarlens}[title={What you've built}]
One more word for this chapter.
\end{grammarlens}

\subsection*{Your turn}
\begin{itemize}
\raggedright
  \item \emph{Say it:} \emph{wecken}
  \item \emph{Say it:} \emph{wecken}, once more
  \item \emph{Say it:} say \emph{wecken}
  \item \emph{Recall:} say the German for to prepare, then the German for to choose, then say \emph{wecken} again
\end{itemize}

\begin{tcolorbox}[breakable,colback=teal!4,colframe=teal!35!black,title={Before you move on}]
What does wecken mean? (\textquotedblleft{}To wake (someone)\textquotedblright{}.) Say it once more.
\end{tcolorbox}
\section[wiegen]{wiegen — to weigh}

\label{lesson:GE-C191-wiegen}



\begin{quote}
Before the new one: say the German for to choose, then the German for to wake.
\end{quote}

\subsection*{You'll want to know: wiegen}
\textbf{wiegen} — \textquotedblleft{}to weigh\textquotedblright{}.

\begin{grammarlens}[title={What you've built}]
One more word for this chapter.
\end{grammarlens}

\subsection*{Your turn}
\begin{itemize}
\raggedright
  \item \emph{Say it:} \emph{wiegen}
  \item \emph{Say it:} \emph{wiegen}, once more
  \item \emph{Say it:} say \emph{wiegen}
  \item \emph{Recall:} say the German for to choose, then the German for to wake, then say \emph{wiegen} again
\end{itemize}

\begin{tcolorbox}[breakable,colback=teal!4,colframe=teal!35!black,title={Before you move on}]
What does wiegen mean? (\textquotedblleft{}To weigh\textquotedblright{}.) Say it once more.
\end{tcolorbox}
\section[zeichnen]{zeichnen — to draw}

\label{lesson:GE-C191-zeichnen}



\begin{quote}
Before the new one: say the German for to wake, then the German for to weigh.
\end{quote}

\subsection*{You'll want to know: zeichnen}
\textbf{zeichnen} — \textquotedblleft{}to draw\textquotedblright{}.

\begin{grammarlens}[title={What you've built}]
One more word for this chapter.
\end{grammarlens}

\subsection*{Your turn}
\begin{itemize}
\raggedright
  \item \emph{Say it:} \emph{zeichnen}
  \item \emph{Say it:} \emph{zeichnen}, once more
  \item \emph{Say it:} say \emph{zeichnen}
  \item \emph{Recall:} say the German for to wake, then the German for to weigh, then say \emph{zeichnen} again
\end{itemize}

\begin{tcolorbox}[breakable,colback=teal!4,colframe=teal!35!black,title={Before you move on}]
What does zeichnen mean? (\textquotedblleft{}To draw\textquotedblright{}.) Say it once more.
\end{tcolorbox}
\section[(dialogue)]{Words from chapters 183-186 — a first pass}

\label{lesson:GE-R191-first-pass-words-from-chapters-183-186}



\begin{quote}
Say the words for to weigh and to draw.
\end{quote}

\subsection*{Your turn}
Say these aloud:

\begin{itemize}
\raggedright
  \item to suffer, to measure, to mix, to sew, to sneeze
  \item to park, to look after, to pick, to test, to clean
  \item to rescue, to row, to push, to hit, to taste
  \item to snow, to protect, to sail, to send, to stroll
\end{itemize}

\begin{tcolorbox}[breakable,colback=teal!4,colframe=teal!35!black,title={Before you move on}]
To suffer? (\textbf{leiden}.) To measure? (\textbf{messen}.) To mix? (\textbf{mischen}.) To sew? (\textbf{nähen}.) To sneeze? (\textbf{niesen}.) To park? (\textbf{parken}.) To look after? (\textbf{pflegen}.) To pick? (\textbf{pflücken}.) To test? (\textbf{prüfen}.) To clean? (\textbf{reinigen}.) To rescue? (\textbf{retten}.) To row? (\textbf{rudern}.) To push? (\textbf{schieben}.) To hit? (\textbf{schlagen}.) To taste? (\textbf{schmecken}.) To snow? (\textbf{schneien}.) To protect? (\textbf{schützen}.) To sail? (\textbf{segeln}.) To send? (\textbf{senden}.) To stroll? (\textbf{spazieren}.)
\end{tcolorbox}
\section[(dialogue)]{Words from chapters 187-191 — a second pass}

\label{lesson:GE-R191-second-pass-words-from-chapters-187-191}



\begin{quote}
Say the words for to rinse and to start.
\end{quote}

\subsection*{Your turn}
Say these aloud:

\begin{itemize}
\raggedright
  \item to rinse, to start, to put, to steal, to climb
  \item to paint, to share, to talk on the phone, to separate, to cross
  \item to translate, to hug, to arrange, to forbid, to connect
  \item to compare, to injure, to miss, to postpone, to use
  \item to prepare, to choose, to wake, to weigh, to draw
\end{itemize}

\begin{tcolorbox}[breakable,colback=teal!4,colframe=teal!35!black,title={Before you move on}]
To rinse? (\textbf{spülen}.) To start? (\textbf{starten}.) To put? (\textbf{stecken}.) To steal? (\textbf{stehlen}.) To climb? (\textbf{steigen}.) To paint? (\textbf{streichen}.) To share? (\textbf{teilen}.) To talk on the phone? (\textbf{telefonieren}.) To separate? (\textbf{trennen}.) To cross? (\textbf{überqueren}.) To translate? (\textbf{übersetzen}.) To hug? (\textbf{umarmen}.) To arrange? (\textbf{verabreden}.) To forbid? (\textbf{verbieten}.) To connect? (\textbf{verbinden}.) To compare? (\textbf{vergleichen}.) To injure? (\textbf{verletzen}.) To miss? (\textbf{verpassen}.) To postpone? (\textbf{verschieben}.) To use? (\textbf{verwenden}.) To prepare? (\textbf{vorbereiten}.) To choose? (\textbf{wählen}.) To wake? (\textbf{wecken}.) To weigh? (\textbf{wiegen}.) To draw? (\textbf{zeichnen}.)
\end{tcolorbox}
